\begin{lemma}[Case 1: the first vertex and equations (13)--(15)]
Use manuscript stages \(j=m(v_n)\leq n\), represented in Lean by
distinct zero-based predecessor slots \(j_{\rm slot}\leq q\), with
\(q.\mathrm{val}+1=n\) and
\(j_{\rm slot}.\mathrm{val}+1=j\).  Let
\[
 L_1=\{1_1<\cdots<1_{p+1}=n\}
\]
be the exact class of active PE slots supported at \(v_n\), with exact
fields \(\gamma_{1_{a+1}}=\gamma_a\),
\(\delta_{1_{a+1}}=\delta_a\), and
\(S_{1_{a+1}}=\{\gamma_a,\delta_a\}\).  Assume
\[
 \gamma_{a+1}=\delta_a\quad(0\leq a<p),
 \qquad \gamma_0,\delta_0,\ldots,\delta_p
             \text{ are pairwise distinct}.                  \tag{4}
\]
Let \(T'_j\), with its displayed closure certificate, be the elementary
closure supplied by \noderef{ExtensionBaseClosure}.  Suppose Case 1 holds:
\[
 \overline\varrho_1(T'_j-v_n)\cap
          \bigcup_{i\in L_1}S_i\ne\varnothing.               \tag{C1}
\]
Assume also the Claim 4.3 missing-color equality on
\(T_j-v_n\), the elementariness of \(T_j\) under
\(\varphi_{j-1}\), the Part I(2) exclusions
\(\delta_a\notin\overline\varphi_{j-1}(T_j)\), and
\[
 \gamma_0\in\overline\varphi_{j-1}(v_n),\qquad
 \overline\varrho_1(v_n)\cap\bigcup_{i\in L_1}S_i
       =\{\delta_p\}.                                        \tag{U}
\]

Then there are a first vertex \(u\) in the order of \(T'_j-v_n\) and
a least position \(b\) such that
\(\gamma_b\in\overline\varrho_1(u)\).  No earlier
\(\gamma_a\), \(a<b\), is missing at \(u\), and every other vertex
of \(T'_j-v_n\) meeting the displayed support-color union is either
\(u\) or occurs after \(u\).  Moreover,
\[
                         u\in V(T'_j)-V(T_j),                 \tag{13}
\]
\[
 \overline\varrho_1(T'_j(u)-u)\cap
  \left(\bigcup_{i\in L_1}S_i-\{\delta_p\}\right)
                         =\varnothing,                         \tag{14}
\]
and, for \(L_1^*=\noderef{ExtensionBaseCaseOneIndices}\),
\[
 \delta_p\notin\bigcup_{i\in L_1^*}S_i,qquad
 \overline\varrho_1(v_n)\cap
       \bigcup_{i\in L_1^*}S_i=\varnothing.                  \tag{15}
\]
\end{lemma}
\begin{proof}
Put \(S(L_1)=\bigcup_{i\in L_1}S_i\).  By (C1), the finite set
\[
 A=\{x\in V(T'_j)\setminus\{v_n\}:
          \overline\varrho_1(x)\cap S(L_1)\ne\varnothing\}
\]
is nonempty.  The duplicate-free vertex list in
\noderef{TreeSequence} gives
a strict total order on its vertex set, as encoded by
\noderef{TreeVertexOrder}.  Therefore \(A\) has a first member \(u\).
This gives \(u\ne v_n\), the asserted minimality property, and the
membership witness \(u\in V(T'_j)\).

Equation (U) puts \(\delta_p\) in
\(\overline\varrho_1(v_n)\).  Since \(V(T'_j)\) is elementary under
\(\varrho_1\), the definition
\noderef{ColoredVertexSetElementary}, applied to the distinct vertices
\(v_n,u\), shows that \(\delta_p\notin
\overline\varrho_1(u)\).  Choose
\(c\in\overline\varrho_1(u)\cap S(L_1)\).  By the exact stored-pair
identities and (4), every color of \(S(L_1)\) is one of
\[
                    \gamma_0,\delta_0,\ldots,\delta_p,
\]
and \(\delta_a=\gamma_{a+1}\) for \(a<p\).  Since \(c\ne\delta_p\),
it follows that \(c=\gamma_b\) for some \(0\leq b\leq p\).  Choose the
least such \(b\).  Then
\(\gamma_b\in\overline\varrho_1(u)\), while no
\(\gamma_a\) with \(a<b\) is missing at \(u\).

We prove (13).  If \(b=0\), then \(\gamma_b=\gamma_0\) is missing at
\(v_n\) under \(\varphi_{j-1}\) by hypothesis.  The vertices of
\(T_j\) are elementary under that coloring, so no other vertex of
\(T_j\) misses \(\gamma_0\).  If \(b>0\), equation (4) gives
\(\gamma_b=\delta_{b-1}\), and the Part I(2) hypothesis says that this
color is missing at no vertex of \(T_j\) under
\(\varphi_{j-1}\).  Thus in either case \(u\ne v_n\) and
\(\gamma_b\) is not missing at \(u\) if \(u\in V(T_j)\), because the
Claim 4.3 agreement hypothesis identifies the two missing-color sets at
every vertex of \(T_j-v_n\).  This contradicts the choice of \(b\), so
\(u\notin V(T_j)\), proving (13).

For (14), take a vertex \(x\) of the prefix \(T'_j(u)-u\).  If
\(x\ne v_n\), then \(x\) occurs strictly before \(u\), so minimality
of \(u\) gives
\(\overline\varrho_1(x)\cap S(L_1)=\varnothing\).  If
\(x=v_n\), equation (U) says that the same intersection is exactly
\(\{\delta_p\}\).  After deleting \(\delta_p\) from the color side,
both cases have empty intersection.  Taking the union over the vertices
in \noderef{TreeVerticesBefore} proves (14).

It remains to prove (15).  By definition,
\(L_1^*\) consists of the positions \(a<b\).  Its stored colors belong
to \(\{\gamma_0,\delta_0,\ldots,\delta_{b-1}\}\).  The injectivity in
(4) makes \(\delta_p\) distinct from every color in this list; hence
\(\delta_p\notin\bigcup_{i\in L_1^*}S_i\).  Intersect the containment
\(\bigcup_{i\in L_1^*}S_i\subseteq S(L_1)\) with (U).  The intersection
can contain only \(\delta_p\), which has just been excluded, so it is
empty.  This is the second assertion of (15), and completes all the
Case 1 selection properties.
\end{proof}
